
\begin{table}[!htbp]
\centering

\caption{Summary of the ventral-view zebrafish larvae microscopy dataset used in this study, including the number of biological fish, annotation records, developmental stage, image resolution, annotated skeletal structures, genotype distribution, and imaging characteristics.}

\label{tab:dataset_summary}

\begin{threeparttable}
\small

\begin{tabularx}{\textwidth}{
    >{\raggedright\arraybackslash}p{5.0cm}
    >{\raggedright\arraybackslash}X
}

\toprule
\rowcolor{HeaderGray}
\textbf{Dataset characteristic} &
\textbf{Description} \\
\midrule

\rowcolor{SectionGray}
\multicolumn{2}{l}{\textbf{Biological specimens}} \\
\addlinespace[2pt]

Species &
Zebrafish (\textit{Danio rerio}) larvae \\

Developmental stage &
10 days post-fertilization (10 dpf) \\

Mutant line &
\textit{col10a1a} \\

Number of individual fish &
34 \\

\addlinespace[4pt]
\rowcolor{SectionGray}
\multicolumn{2}{l}{\textbf{Microscopy dataset}} \\
\addlinespace[2pt]

Total number of microscopy images &
190 \\

Images per fish &
4--6 \\

Microscopy modality &
Brightfield microscopy \\

Anatomical orientation &
Ventral craniofacial view \\

Original image resolution &
$2576 \times 1932$ pixels \\

Annotated skeletal structures &
25 segmentation classes \\

\addlinespace[4pt]
\rowcolor{SectionGray}
\multicolumn{2}{l}{\textbf{Genotype distribution}} \\
\addlinespace[2pt]

Wild type (WT) &
50 images \\

Heterozygous (HET) &
110 images \\

Homozygous (HOM) &
32 images \\

\bottomrule

\end{tabularx}

\begin{tablenotes}[flushleft]
\footnotesize
\item dpf, days post-fertilization; WT, wild type; HET, heterozygous; HOM, homozygous.
\end{tablenotes}

\end{threeparttable}
\end{table}